\documentclass{article}

\usepackage{amsmath,amssymb}
\usepackage{xurl}
\usepackage[hidelinks]{hyperref}
\setlength{\emergencystretch}{2em}

% Shared commands for the notes in docs/paper-gaps/.
% Every note loads this file via % Shared commands for the notes in docs/paper-gaps/.
% Every note loads this file via % Shared commands for the notes in docs/paper-gaps/.
% Every note loads this file via \input{command}. Project-wide notation
% belongs here, never in per-note redefinitions.

% --- Project configuration ------------------------------------------------
\newcommand{\projectIssueBase}{https://github.com/LionSR/QICLean/issues/}
\newcommand{\projectPrBase}{https://github.com/LionSR/QICLean/pull/}
\newcommand{\projectDocBase}{https://sirui-lu.com/QICLean/docs/}
% ---------------------------------------------------------------------------

\DeclareUnicodeCharacter{2080}{\ifmmode{}_{0}\else\textsubscript{0}\fi}
\DeclareUnicodeCharacter{2081}{\ifmmode{}_{1}\else\textsubscript{1}\fi}
\DeclareUnicodeCharacter{2082}{\ifmmode{}_{2}\else\textsubscript{2}\fi}
\DeclareUnicodeCharacter{2099}{\ifmmode{}_{n}\else\textsubscript{n}\fi}

\newcommand{\Lean}{\textsc{Lean}}

% A formal declaration name, upright sans, breakable at underscores.
\ExplSyntaxOn
\NewDocumentCommand{\leanid}{m}{
  \ifmmode
    \textsf{\detokenize{#1}}
  \else
    \group_begin:
    \urlstyle{sf}
    \tl_set:Nn \l_tmpa_tl {#1}
    \tl_replace_all:Nnn \l_tmpa_tl {\_} {_}
    \exp_args:NV \path \l_tmpa_tl
    \group_end:
  \fi
}
\ExplSyntaxOff

% Traceability links; use in footnotes, not in the running prose.
\newcommand{\ghissue}[1]{\href{\projectIssueBase#1}{GitHub issue \##1}}
\newcommand{\ghpr}[1]{\href{\projectPrBase#1}{PR \##1}}
\newcommand{\doclink}[1]{\href{\projectDocBase#1.html}{\path{#1}}}

% Domain notation, using \providecommand so a note-local refinement stays
% possible.
\providecommand{\tr}{\operatorname{tr}}

% Dirac notation and the identity operator, as in blueprint/src/macros.
\usepackage{dsfont}
\providecommand{\ket}[1]{|#1\rangle}
\providecommand{\bra}[1]{\langle #1|}
\providecommand{\braket}[2]{\langle #1 | #2 \rangle}
\providecommand{\ketbra}[2]{|#1\rangle\!\langle #2|}
\providecommand{\Id}{\mathds{1}}
\providecommand{\one}{\mathds{1}}
\providecommand{\C}{\mathbb{C}}
\providecommand{\MN}[1]{M_{#1}(\C)}
\providecommand{\MD}{M_D(\C)}

% Verdict marker: \gapnote{<kind>}{<status>}. Kind is the policy
% classification (clarification, local-correction, scope-restriction,
% unfaithful, false-source, open-gap); status is open, wip (elimination
% underway), resolved, or historical. Machine-read by the site index;
% prints a verdict line.
\newcommand{\gapnote}[2]{\par\noindent\textbf{Verdict:} #1 (#2).\par}
. Project-wide notation
% belongs here, never in per-note redefinitions.

% --- Project configuration ------------------------------------------------
\newcommand{\projectIssueBase}{https://github.com/LionSR/QICLean/issues/}
\newcommand{\projectPrBase}{https://github.com/LionSR/QICLean/pull/}
\newcommand{\projectDocBase}{https://sirui-lu.com/QICLean/docs/}
% ---------------------------------------------------------------------------

\DeclareUnicodeCharacter{2080}{\ifmmode{}_{0}\else\textsubscript{0}\fi}
\DeclareUnicodeCharacter{2081}{\ifmmode{}_{1}\else\textsubscript{1}\fi}
\DeclareUnicodeCharacter{2082}{\ifmmode{}_{2}\else\textsubscript{2}\fi}
\DeclareUnicodeCharacter{2099}{\ifmmode{}_{n}\else\textsubscript{n}\fi}

\newcommand{\Lean}{\textsc{Lean}}

% A formal declaration name, upright sans, breakable at underscores.
\ExplSyntaxOn
\NewDocumentCommand{\leanid}{m}{
  \ifmmode
    \textsf{\detokenize{#1}}
  \else
    \group_begin:
    \urlstyle{sf}
    \tl_set:Nn \l_tmpa_tl {#1}
    \tl_replace_all:Nnn \l_tmpa_tl {\_} {_}
    \exp_args:NV \path \l_tmpa_tl
    \group_end:
  \fi
}
\ExplSyntaxOff

% Traceability links; use in footnotes, not in the running prose.
\newcommand{\ghissue}[1]{\href{\projectIssueBase#1}{GitHub issue \##1}}
\newcommand{\ghpr}[1]{\href{\projectPrBase#1}{PR \##1}}
\newcommand{\doclink}[1]{\href{\projectDocBase#1.html}{\path{#1}}}

% Domain notation, using \providecommand so a note-local refinement stays
% possible.
\providecommand{\tr}{\operatorname{tr}}

% Dirac notation and the identity operator, as in blueprint/src/macros.
\usepackage{dsfont}
\providecommand{\ket}[1]{|#1\rangle}
\providecommand{\bra}[1]{\langle #1|}
\providecommand{\braket}[2]{\langle #1 | #2 \rangle}
\providecommand{\ketbra}[2]{|#1\rangle\!\langle #2|}
\providecommand{\Id}{\mathds{1}}
\providecommand{\one}{\mathds{1}}
\providecommand{\C}{\mathbb{C}}
\providecommand{\MN}[1]{M_{#1}(\C)}
\providecommand{\MD}{M_D(\C)}

% Verdict marker: \gapnote{<kind>}{<status>}. Kind is the policy
% classification (clarification, local-correction, scope-restriction,
% unfaithful, false-source, open-gap); status is open, wip (elimination
% underway), resolved, or historical. Machine-read by the site index;
% prints a verdict line.
\newcommand{\gapnote}[2]{\par\noindent\textbf{Verdict:} #1 (#2).\par}
. Project-wide notation
% belongs here, never in per-note redefinitions.

% --- Project configuration ------------------------------------------------
\newcommand{\projectIssueBase}{https://github.com/LionSR/QICLean/issues/}
\newcommand{\projectPrBase}{https://github.com/LionSR/QICLean/pull/}
\newcommand{\projectDocBase}{https://sirui-lu.com/QICLean/docs/}
% ---------------------------------------------------------------------------

\DeclareUnicodeCharacter{2080}{\ifmmode{}_{0}\else\textsubscript{0}\fi}
\DeclareUnicodeCharacter{2081}{\ifmmode{}_{1}\else\textsubscript{1}\fi}
\DeclareUnicodeCharacter{2082}{\ifmmode{}_{2}\else\textsubscript{2}\fi}
\DeclareUnicodeCharacter{2099}{\ifmmode{}_{n}\else\textsubscript{n}\fi}

\newcommand{\Lean}{\textsc{Lean}}

% A formal declaration name, upright sans, breakable at underscores.
\ExplSyntaxOn
\NewDocumentCommand{\leanid}{m}{
  \ifmmode
    \textsf{\detokenize{#1}}
  \else
    \group_begin:
    \urlstyle{sf}
    \tl_set:Nn \l_tmpa_tl {#1}
    \tl_replace_all:Nnn \l_tmpa_tl {\_} {_}
    \exp_args:NV \path \l_tmpa_tl
    \group_end:
  \fi
}
\ExplSyntaxOff

% Traceability links; use in footnotes, not in the running prose.
\newcommand{\ghissue}[1]{\href{\projectIssueBase#1}{GitHub issue \##1}}
\newcommand{\ghpr}[1]{\href{\projectPrBase#1}{PR \##1}}
\newcommand{\doclink}[1]{\href{\projectDocBase#1.html}{\path{#1}}}

% Domain notation, using \providecommand so a note-local refinement stays
% possible.
\providecommand{\tr}{\operatorname{tr}}

% Dirac notation and the identity operator, as in blueprint/src/macros.
\usepackage{dsfont}
\providecommand{\ket}[1]{|#1\rangle}
\providecommand{\bra}[1]{\langle #1|}
\providecommand{\braket}[2]{\langle #1 | #2 \rangle}
\providecommand{\ketbra}[2]{|#1\rangle\!\langle #2|}
\providecommand{\Id}{\mathds{1}}
\providecommand{\one}{\mathds{1}}
\providecommand{\C}{\mathbb{C}}
\providecommand{\MN}[1]{M_{#1}(\C)}
\providecommand{\MD}{M_D(\C)}

% Verdict marker: \gapnote{<kind>}{<status>}. Kind is the policy
% classification (clarification, local-correction, scope-restriction,
% unfaithful, false-source, open-gap); status is open, wip (elimination
% underway), resolved, or historical. Machine-read by the site index;
% prints a verdict line.
\newcommand{\gapnote}[2]{\par\noindent\textbf{Verdict:} #1 (#2).\par}


\title{The Commutation Hypothesis in Wolf's Trace Jensen Inequality}
\author{}
\date{2026-11-15}

\begin{document}
\maketitle
\gapnote{false-source}{resolved}

\section{The assertion}

Let $T_1,\ldots,T_r:\MN{d}\to\MN{d'}$ be positive linear maps, let
$f:I\to\mathbb R$ be convex on an interval $I\subseteq\mathbb R$, and let
$X_1,\ldots,X_r\in\MN{d}$ be Hermitian with spectra in $I$. Write
\begin{align}
    Y&=\sum_{i=1}^rT_i(X_i).
    \label{eq:wolf_ch5_trace_jensen_commutation_argument}
\end{align}
Wolf's Theorem~5.17 asserts
\begin{align}
    \tr\bigl[C f(Y)\bigr]
        &\leq\sum_{i=1}^r\tr\bigl[C\,T_i(f(X_i))\bigr]
    \label{eq:wolf_ch5_trace_jensen_commutation_printed_inequality}
\end{align}
for every positive semidefinite $C\in\MN{d'}$ which commutes with $f(Y)$,
provided either $\sum_iT_i(\Id)=\Id$, or $\sum_iT_i(\Id)\leq\Id$ and
$f(0)\leq0$. This is Equation~(5.51) of Ref.~\cite{Wolf2012Quantum}.\footnote{Source:
    \cite[Theorem~5.17 and Equations~(5.51)--(5.53)]{Wolf2012Quantum}; local
    source \path{Notes/WolfNoteTexSource/ch05_schwarz_inequalities.tex},
    lines~1039--1078. Tracked in \ghissue{27}.}

\section{The point at issue}

The proof reduces to $C=\ketbra{\psi}{\psi}$ for a unit eigenvector $\psi$ of
$f(Y)$ and then uses
\begin{align}
    \bra{\psi}f(Y)\ket{\psi}
        &=f\bigl(\bra{\psi}Y\ket{\psi}\bigr).
    \label{eq:wolf_ch5_trace_jensen_commutation_eigenvector_step}
\end{align}
This is Equation~(5.52) of Ref.~\cite{Wolf2012Quantum}. It holds when $\psi$
is an eigenvector of $Y$. An eigenvector of $f(Y)$ need not be one when $f$
takes the same value at two eigenvalues of $Y$, and then the printed
inequality fails.

Take $d=1$, $d'=2$, $r=3$, $I=[-2,2]$, $f(x)=|x|$, and the scalars
$X_1=2$, $X_2=-2$, $X_3=0$. Define $T_i(x)=xP_i$, where
\begin{align}
    P_1&=\frac1{16}\begin{pmatrix}9&2\\2&1\end{pmatrix},
    &
    P_2&=\frac1{16}\begin{pmatrix}1&2\\2&9\end{pmatrix},
    &
    P_3&=\frac1{16}\begin{pmatrix}6&-4\\-4&6\end{pmatrix}.
    \label{eq:wolf_ch5_trace_jensen_commutation_weights}
\end{align}
Each $P_i$ is positive semidefinite, so each $T_i$ is positive, and
$\sum_iT_i(1)=P_1+P_2+P_3=\Id$. Moreover
\begin{align}
    Y&=2P_1-2P_2=\begin{pmatrix}1&0\\0&-1\end{pmatrix},
    &
    f(Y)&=\Id.
    \label{eq:wolf_ch5_trace_jensen_commutation_identity_image}
\end{align}
Hence every matrix commutes with $f(Y)$. For the positive semidefinite
matrix
\begin{align}
    C&=\begin{pmatrix}1&-1\\-1&1\end{pmatrix}
    \label{eq:wolf_ch5_trace_jensen_commutation_test_matrix}
\end{align}
we obtain
\begin{align}
    \sum_{i=1}^3\tr\bigl[C\,T_i(f(X_i))\bigr]
        &=\tr\Bigl[C\,\frac1{16}\begin{pmatrix}20&8\\8&20\end{pmatrix}\Bigr]
        =\frac32
        <2=\tr\bigl[Cf(Y)\bigr].
    \label{eq:wolf_ch5_trace_jensen_commutation_violation}
\end{align}
All hypotheses of the unital alternative hold, but
(\ref{eq:wolf_ch5_trace_jensen_commutation_printed_inequality}) fails.

\section{The corrected statement and source proof}

The source argument proves
(\ref{eq:wolf_ch5_trace_jensen_commutation_printed_inequality}) for every
positive semidefinite $C$ commuting with $Y$ itself. Such a $C$ commutes with
$f(Y)$, and the two conditions agree when $f$ is injective on $I$. In the
subunital alternative the source evaluates $f(0)$, so $0\in I$ is part of
that alternative.

Write $C=S^2$ with $S=C^{1/2}$. Then $S$ commutes with $Y$ and with $f(Y)$.
Let $e_1,\ldots,e_{d'}$ be an orthonormal eigenbasis of $Y$ with
$Ye_k=y_ke_k$, and set $v_k=Se_k$. Then $Yv_k=y_kv_k$ and
$f(Y)v_k=f(y_k)v_k$, so
\begin{align}
    \tr\bigl[Cf(Y)\bigr]
        &=\sum_k\bra{v_k}f(Y)\ket{v_k}
        =\sum_k\lVert v_k\rVert^2f(y_k),
    \label{eq:wolf_ch5_trace_jensen_commutation_left_side}\\
    \sum_i\tr\bigl[C\,T_i(f(X_i))\bigr]
        &=\sum_k\sum_i\bra{v_k}T_i(f(X_i))\ket{v_k}.
    \label{eq:wolf_ch5_trace_jensen_commutation_right_side}
\end{align}
Write $X_i=\sum_x\mu_{x,i}Q_{x,i}$ with rank-one spectral projections and set
$p_{x,i}=\bra{v_k}T_i(Q_{x,i})\ket{v_k}\geq0$. Then
$\sum_{x,i}p_{x,i}\mu_{x,i}=\lVert v_k\rVert^2y_k$ and
$\sum_{x,i}p_{x,i}=\bra{v_k}\sum_iT_i(\Id)\ket{v_k}$, which equals
$\lVert v_k\rVert^2$ or is at most $\lVert v_k\rVert^2$ in the two
alternatives. Scalar Jensen, with the missing mass placed at $0\in I$ in the
subunital case, gives
\begin{align}
    \lVert v_k\rVert^2f(y_k)
        &\leq\sum_{x,i}p_{x,i}f(\mu_{x,i})
        =\sum_i\bra{v_k}T_i(f(X_i))\ket{v_k}.
    \label{eq:wolf_ch5_trace_jensen_commutation_pointwise}
\end{align}
Summing over $k$ proves the corrected inequality.

\section{Formalization}

The corrected theorem is
\leanid{Matrix.re\_trace\_mul\_cfc\_sum\_map\_le\_of\_commute}, and the
counterexample is
\leanid{Matrix.trace\_jensen\_commute\_cfc\_counterexample}.\footnote{Formal
    module: \doclink{QICLean/Channel/Schwarz/TraceJensen}.}
The scalar steps are \leanid{ConvexOn.mul\_map\_le\_sum\_of\_sum\_eq} and
\leanid{ConvexOn.mul\_map\_le\_sum\_of\_sum\_le}. The finite family of maps is
indexed by a finite type, the traces are compared through their real parts,
and no further hypothesis is added.

\section{Conclusion}

Wolf's Theorem~5.17 is false as printed: commutation of $C$ with $f(Y)$ does
not suffice. Requiring that $C$ commute with $Y$, which is what
Equation~(5.52) uses, makes the source proof valid, and this corrected
theorem is formalized.

\bibliographystyle{alpha}
\bibliography{references}

\end{document}
